\documentclass[a4paper, 12pt]{article}

%%% BuS LAB PROTOCOLL PREAMBLE
%%% 2019
%%%%%%%%%%%%%%%%%%%%%%%%%%%%%%%


%%% PACKAGES
%%%%%%%%%%%%%%%%%%%%%%%%%%%

\usepackage[ngerman]{babel}

\usepackage[utf8]{inputenc}
\usepackage{amsmath}
\usepackage{pgfplots}
\usepackage{tikz}
\usepackage[many]{tcolorbox}
\usepackage{graphicx,import}
\graphicspath{ {./graphics/} }
\usepackage{pdfpages}
\usepackage{dashrule}
\usepackage{float}
\usepackage{siunitx}
\usepackage{booktabs}
\usepackage{transparent}
\usepackage{svg}
\usepackage[version=4]{mhchem}
\usepackage[bottom]{footmisc}
\usepackage[european, straightvoltages]{circuitikz}
\usepackage{multicol}


%%% DOCUMENT GEOMETRY
%%%%%%%%%%%%%%%%%%%%%%%%%%%

\usepackage{geometry}
\geometry{
 a4paper,
 total={0.6180339887498948\paperwidth,0.6180339887498948\paperheight},
 top = 0.1458980337503154\paperheight,
 bottom = 0.1458980337503154\paperheight
 }
\setlength{\jot}{0.013155617496424828\paperheight}
\linespread{1.1458980337503154}

\setlength{\parskip}{0.013155617496424828\paperheight} % paragraph spacing


%%% COLORS
%%%%%%%%%%%%%%%%%%%%%%%%%%%

\definecolor{red1}{HTML}{f38181}
\definecolor{yellow1}{HTML}{fce38a}
\definecolor{green1}{HTML}{95e1d3}
\definecolor{blue1}{HTML}{66bfbf}
\definecolor{hsblue}{HTML}{00b1db}
\definecolor{hsgrey}{HTML}{afafaf}

%%% CONSTANTS
%%%%%%%%%%%%%%%%%%%%%%%%%%%
\newlength{\smallvert}
\setlength{\smallvert}{0.0131556\paperheight}


%%% COMMANDS
%%%%%%%%%%%%%%%%%%%%%%%%%%%

% differential d
\newcommand*\dif{\mathop{}\!\mathrm{d}}

% horizontal line
\newcommand{\holine}[1]{
  	\begin{center}
	  	\noindent{\color{hsgrey}\hdashrule[0ex]{#1}{1pt}{3mm}}\\%[0.0131556\paperheight]
  	\end{center}
}

% mini section
\newcommand{\minisec}[1]{ \noindent\underline{\textit {#1} } \\}

% quick function plot
\newcommand{\plotfun}[3]{
  \vspace{0.021286\paperheight}
  \begin{center}
    \begin{tikzpicture}
      \begin{axis}[
        axis x line=center,
        axis y line=center,
        ]
        \addplot[draw=red1][domain=#2:#3]{#1};
      \end{axis}
    \end{tikzpicture}
  \end{center}
}

% box for notes
\newcommand{\notebox}[1]{

\tcbset{colback=white,colframe=red1!100!black,title=Note!,width=0.618\paperwidth,arc=0pt}

 \begin{center}
  \begin{tcolorbox}[]
   #1 
  \end{tcolorbox}
 
 \end{center} 
 
}

% box for equation
\newcommand{\eqbox}[2]{
	
	\tcbset{colback=white,colframe=hsblue!100!black,title=,width=#2,arc=0pt}
	
	\begin{center}
		\begin{tcolorbox}[ams align*]
				#1
		\end{tcolorbox}
		
	\end{center} 
	
}

\newcommand{\inklude}[1]{%
    \def\svgwidth{\columnwidth}
    \import{./graphics/}{#1.pdf_tex}
}
% END OF PREAMBLE

%%%%%%%%%%%%%%%%%%%%%%%%%%%%%%%%%%%%%

\begin{document}

%%%%%%%%%%%%%%%%%%%%%%%%%%%%%%%%%%%%%
  \includepdf{./titlepage/titlepage.pdf}
  \clearpage
  \setcounter{page}{1}
%%%%%%%%%%%%%%%%%%%%%%%%%%%%%%%%%%%%%


  \section{Register für die Timergesteuerte AD-Wandlung}

  \begin{itemize}

  \item \inlinecode{P6DIR} \commentGray{(Input Direction)}
  \item \inlinecode{P6SEL} \commentGray{(Alternativfunktion: ADC)}

  \item \inlinecode{ADC12CTL0}
    \begin{itemize}
      \item \inlinecode{SHT0x, SHT1x}
      \item \inlinecode{REF2_5V}
      \item \inlinecode{REFON}
      \item \inlinecode{ADC12ON}
      \item \inlinecode{ENC} \commentGray{Konfiguration}
      \item \inlinecode{ADC12SC}
    \end{itemize}

  \item \inlinecode{ADC12CTL1}
    \begin{itemize}
      \item \inlinecode{SHP}
      \item \inlinecode{SHSx}
      \item \inlinecode{ADC12SSELx}
      \item \inlinecode{CONSEQx} \commentGray{Modusauswahl}
      \item \inlinecode{CSTARTADDx} 
    \end{itemize}

  \item \inlinecode{ADC12MCTLx} \commentGray{(conversion memory control register(s))}
    \begin{itemize}
    \item \inlinecode{SREF0, SREF1, SREF2}
    \item \inlinecode{INCHx} \commentGray{(P6.3)}
    \item \inlinecode{EOS}
    \end{itemize}

  \item \inlinecode{ADC12IE}
  \item \inlinecode{TACTL} \commentGray{Abtastperiodeneinstellung}

  \end{itemize}

  \section{Umrechnung des AD-Wertes in eine Spannung}
  Die Ermittlung des Spannungswertes eines gewandelten digitalen Wertes gelingt durch
  \[ U = AD_{d} \cdot \mathrm{LSB} = AD_{d} \cdot \frac{\mathrm{FSR}}{2^{12}}\]

  Bei z.B. einer FSR von $5 \, \si{\volt}$ und einem gewandelten Wert von
  $AD_{d} = 4095_d = FFF_h $ erhält man eine Spannung von $U=4.9988 \, \si{\volt}$

  \section{Einlesbarer Spannungsbereich}
  Bei Wandlung mit einem SAR-Wandler, an dessen DAC eine Referenzspannung von
  $U_{\mathrm{ref}} = 2.5 \, \si{\volt}$ liegt, ist der gültig einlesbare
  Spannungsbereich auf diesen Referenzwert begrenzt, da der DAC eine maximale
  Spannung von $U_{\mathrm{ref}}$ an den Komparator rückkoppeln kann und der
  konvertierte Wert bei einer Eingangsspannung größer $U_\mathrm{ref}$ dann nur
  zu einem maximalen Wert \inlinecode{111111111111} gewandelt wird, welcher
  $U_{\mathrm{ref}}$ repräsentiert.

  \end{document}

% Local Variables:
% TeX-engine: luatex
% End:
